Этот раздел даёт компактный содержательный обзор постановки и ключевых
обозначений; формальные определения и детерминированные предпосылки вводятся в
разделе \emph{Formal Model}. Модель описывает фиксированный объём программной работы, доводимый до
заранее заданного критерия приёмки. Исходным сценарием служит последовательное
выполнение того же объёма одним разработчиком без ИИ. Полный словарь терминов
и развёрнутая система обозначений приведены в
Приложении~\ref{app:glossary}.

\begin{description}[style=nextline]
  \item[Задача и режим.] Задача $i$ имеет сложность $Z_i$ и выполняется в
  режиме $r$: интерактивном, делегированном либо полностью автономном. Режим
  задаёт порядок постановки, агентного исполнения, проверки и исправления
  результата. Основной анализ M4 относится к интерактивному и делегированному
  режимам; полностью автономный режим без обязательной human-фазы служит
  предельным случаем.

  \item[Нормированная длительность с ИИ.] Величина
  $k_i^{(r)}=T_i^{ai}(r)/(Z_iX)$ сравнивает длительность принятого результата
  с базовым временем той же задачи. Она относится к тройке «задача --
  инструмент -- режим», а не к ИИ-инструменту вообще. Это описательная
  нормировка; причинная интерпретация требует отдельного дизайна сравнения.

  \item[Фазы.] Нормированная длительность раскладывается как
  $k_i^{(r)}=a_i^{(r)}+h_i^{(r)}$. Автономные agent-фазы требуют агентного
  слота, human-фазы --- одновременно внимания разработчика и сохранения
  назначенного слота.

  \item[Управляемый параллелизм.] Параметр $P$ обозначает число агентных
  потоков, настроенных в данной конфигурации и доступных расписанию; для этого
  значения оцениваются координационные издержки. Неиспользуемая дополнительная
  мощность при неизменной конфигурации сама по себе overhead не создаёт.

  \item[Конфигурация.] Запись $\sigma=(P,r)$ или $\sigma=(P,\rho)$ является
  сокращённой при фиксированных ИИ-инструменте и правилах приёмки. Если
  инструмент или критерий принятого результата меняются, их необходимо явно
  указывать при сравнении, хотя они не добавляются в размерность $\sigma$.

  \item[Расписание и срок.] Допустимое расписание соблюдает зависимости DAG,
  порядок фаз, мощность человеческого ресурса 1 и локальное удержание слота.
  Его срок обозначается $T(\pi)$, а минимальный срок среди допустимых
  расписаний --- $T^*$.
\end{description}

\begin{center}
\small
\begin{tabularx}{\textwidth}{@{}>{$}l<{$}>{\raggedright\arraybackslash}X@{}}
\toprule
\text{Обозначение} & \text{Смысл} \\
\midrule
X & базовая трудоёмкость единицы работы без ИИ \\
Z_i & сложность задачи $i$ в единицах $X$ \\
k_i=a_i+h_i & нормированная длительность и её agent/human-составляющие \\
P & управляемое число агентных потоков конфигурации \\
G=(V,E) & ориентированный ациклический граф задач \\
W=X\sum_iZ_ik_i & суммарная работа при $C=\gamma=1$ \\
A,\ H & суммарная автономная агентная и человеческая работа \\
C(P),\ \gamma(P) & множители межагентных и человеческих издержек \\
T(\pi),\ T^* & срок конкретного расписания и оптимальный makespan \\
\bottomrule
\end{tabularx}
\end{center}
